% !TEX root = ../../tree_notes.tex
% =====================================================================
%  Chapter 1 opener -- "Chapter at a Glance", reading paths + ladder
%  staircase, roadmap, topic table, recognition guide, notation.
%  NO numbered boxes here (they would be numbered 1.0.x).
% =====================================================================

\begin{chaptertheorybox}{Chapter at a Glance}
A \term{tree} is the simplest structure that captures a hierarchy: one thing at the
top, everything else hanging below it, and exactly one way to travel from any object
to any other. This chapter builds the whole vocabulary we use for the rest of the
book -- \emph{root, parent, child, leaf, subtree, ancestor, descendant, depth,
height, size} -- and proves the few facts that every later algorithm quietly relies
on: a tree with $n$ nodes has $n-1$ edges, there is exactly one path between any two
nodes, and height can be computed bottom-up with a single recursion.

\textbf{You will learn to}
\begin{itemize}[nosep]
  \item say precisely which pictures are trees and which are not, and why;      \hfill\lv{1}
  \item prove and use the $n-1$ edges and unique-path properties;               \hfill\lv{1}
  \item compute $\depth$, $\height$ and $\size$ of every node in one DFS;        \hfill\lv{2}
  \item name the standard shapes (binary, full, complete, perfect, balanced, skewed)
        and know the counting rule each one gives;                    \hfill\lv{2}--\lv{3}
  \item count tree shapes with the Catalan numbers and labelled trees with Cayley's
        $n^{\,n-2}$ formula;                                                    \hfill\lv{3}--\lv{4}
  \item recognise a \emph{tree-in-disguise} statement and route it to the right tool. \hfill\lv{4}
\end{itemize}

\textbf{Prerequisites.} None -- this is the entry point of the book. We only assume
you can read a small $\cpp$ function and know what a graph (nodes and edges) is.

\textbf{In contests and interviews.} Every tree problem on LeetCode and most on
Codeforces start here: you must know what a root, a leaf and a subtree are, and that
a DFS carrying the parent is the workhorse. The counting facts of \secref{c01-counting}
turn up directly in combinatorics and DP problems (\chref{c10}).
\end{chaptertheorybox}

% ---------------- Station 1b: reading paths + ladder position ----------------
\begin{chaptertheorybox}{How to Read This Chapter (two paths up the ladder)}
Every problem, solution step and warning in this book carries a level badge
$\lv1$--$\lv6$: $\lv1$ is \emph{``what is a tree''} (rating about 800), $\lv3$ is the
solid interview medium (about 1600), $\lv5$--$\lv6$ is Codeforces 2600--3500+.

\textbf{Beginner path ($\lv1$--$\lv3$).} Read \secref{c01-what-is-tree} and
\secref{c01-anatomy} completely; in \secref{c01-kinds} read the definitions and the
first two solved problems; in \secref{c01-counting} you may \emph{skip on the first
pass}: the Catalan recurrence proof, Cayley's proof, the $\lv4$+ variation entries,
the hardening notes and the chapter's Summit. Do every recall card and the
$\lv1$--$\lv3$ practice items. At the end of this path you can solve every
$\lv3$ tree-recursion problem in this chapter.

\textbf{Full path.} Everything, in order -- it ends with a scaled summit
(the hardest problem solvable with Chapter-1 tools alone) and the reduction gallery
that shows how often ``it is a tree'' is the whole trick. Chapter~1 tops out at
$\lv4$--$\lv5$; the true $\lv6$ summits begin once Chapters 13--16 hand you DP,
lifting and query trees.
\end{chaptertheorybox}

\begin{figure}[H]
\centering
\begin{tikzpicture}[diagram, x=9mm, y=7mm]
  \foreach \st/\lab in {1/lv1,2/lv2,3/lv3,4/lv4,5/lv5,6/lv6} {
    \fill[\lab!22, draw=\lab, line width=0.8pt] (\st-1,0) rectangle (\st,\st*0.52);
    \node[font=\scriptsize\bfseries, text=\lab] at (\st-0.5,\st*0.52-0.42) {L\st};}
  \node[annot, anchor=south west, font=\scriptsize]
    at (0.05,0.62) {1.1 what is};
  \node[annot, anchor=south west, font=\scriptsize]
    at (1.05,1.14) {1.2 anatomy};
  \node[annot, anchor=south west, font=\scriptsize]
    at (2.05,1.66) {1.3 kinds};
  \node[annot, anchor=south west, font=\scriptsize]
    at (3.05,2.18) {1.4 counting};
  \node[annot, anchor=south west, font=\scriptsize, text=lv5]
    at (4.05,2.70) {scaled summit};
  \node[annot, anchor=south west, font=\scriptsize, text=lv6]
    at (5.05,3.22) {true summit (later chs)};
\end{tikzpicture}
\caption*{Chapter~1 on the six-step ladder: it teaches steps $\lv1$--$\lv4$ and
shows $\lv5$ from the bottom. Beginner path $=$ steps 1--3; the scaled summit sits
on step~5 and waits for you after the full path.}
\end{figure}

% ---------------- Roadmap diagram ----------------
\begin{figure}[H]
\centering
\begin{tikzpicture}[diagram, node distance=4mm,
    top/.style={boxnode, fill=primaryblue!10, draw=primaryblue, text width=25mm,
                minimum height=9mm, font=\small},
    res/.style={annot, text width=25mm, anchor=north}]
  \node[top] (t1) {\textbf{1.1} What is a tree?};
  \node[top, right=of t1] (t2) {\textbf{1.2} Depth, height, size};
  \node[top, right=of t2] (t3) {\textbf{1.3} Kinds of trees};
  \node[top, right=of t3] (t4) {\textbf{1.4} Counting trees};
  \foreach \a/\b in {t1/t2, t2/t3, t3/t4} \draw[flow] (\a) -- (\b);
  \node[res] at ([yshift=-2mm]t1.south) {$n$ nodes, $n-1$ edges; unique path;
    rooted vs.\ unrooted};
  \node[res] at ([yshift=-2mm]t2.south) {depth from root, height from leaves,
    bottom-up recursion};
  \node[res] at ([yshift=-2mm]t3.south) {binary, full, complete, perfect,
    balanced, skewed};
  \node[res] at ([yshift=-2mm]t4.south) {Catalan $C_n$ shapes; Cayley $n^{n-2}$
    labelled trees};
\end{tikzpicture}
\caption*{The four topics of this chapter, left to right, each with the results you
take away from it.}
\end{figure}

% ---------------- Topic table ----------------
\begin{center}
\renewcommand{\arraystretch}{1.25}
\begin{tabular}{@{}p{0.26\linewidth} p{0.36\linewidth} c p{0.22\linewidth}@{}}
\toprule
\textbf{Topic} & \textbf{Key results} & \textbf{Levels} & \textbf{Solved problems} \\
\midrule
\secref{c01-what-is-tree}\ \ What is a tree?
  & $n-1$ edges; exactly one path between any two nodes
  & $\lv1$--$\lv2$
  & \probref{c01-same-tree} \\
\addlinespace
\secref{c01-anatomy}\ \ Depth, height and size
  & height/size recurrences; one-DFS measurement
  & $\lv1$--$\lv3$
  & \probref{c01-max-depth}, \probref{c01-min-depth} \\
\addlinespace
\secref{c01-kinds}\ \ Kinds of trees
  & perfect $2^{h+1}-1$ nodes; full tree leaves $=$ internal $+1$
  & $\lv2$--$\lv4$
  & \probref{c01-balanced}, \probref{c01-count-nodes} \\
\addlinespace
\secref{c01-counting}\ \ Counting trees
  & Catalan $C_n$ shapes; Cayley $n^{n-2}$ labelled trees
  & $\lv3$--$\lv5$
  & \probref{c01-unique-bst} \\
\bottomrule
\end{tabular}
\end{center}

% ---------------- Problem recognition guide ----------------
\subsection*{Problem Recognition Guide --- When You See \emph{this}, Think \emph{that}}

The fastest way to get stuck on a tree problem is to reach for the wrong tool. Before
you code, match the problem's phrasing against this table; every row is a pattern you
will formally build in the section named at the right.

\begin{center}
\renewcommand{\arraystretch}{1.3}
\begin{tabular}{@{}p{0.40\linewidth} p{0.38\linewidth} l@{}}
\toprule
\textbf{When the problem says\ldots} & \textbf{\textcolor{deepblue}{\textbf{Think\ldots}}} & \textbf{Section} \\
\midrule
``same / identical / duplicate subtree'', compare two trees
  & structural recursion walking both trees in lockstep & \secref{c01-what-is-tree} \\
\addlinespace
``for each node, how far from the root / which level''
  & carry \code{depth + 1} \emph{down} the recursion (top-down) & \secref{c01-anatomy} \\
\addlinespace
``height / deepest / how many nodes below / subtree''
  & \emph{return} a value \emph{up} the recursion (bottom-up) & \secref{c01-anatomy} \\
\addlinespace
``the tree is complete / perfect / full / balanced''
  & the shape's counting rule ($2^{h+1}{-}1$, $L{=}I{+}1$) skips work & \secref{c01-kinds} \\
\addlinespace
``how many different tree shapes / BSTs from $1..n$ / labelled trees''
  & Catalan $C_n$ or Cayley $n^{\,n-2}$, usually with a modulus & \secref{c01-counting} \\
\addlinespace
$n \le 10^{5}$\textasciitilde$10^{6}$ and the tree may be a long chain
  & recursion depth is the danger -- iterative DFS or an explicit stack & \secref{c01-anatomy} \\
\bottomrule
\end{tabular}
\end{center}

\begin{center}
\renewcommand{\arraystretch}{1.25}
\begin{tabular}{@{}p{0.26\linewidth} p{0.46\linewidth} p{0.20\linewidth}@{}}
\toprule
\textbf{Constraint you see} & \textbf{What it permits} & \textbf{Chapter-1 tool} \\
\midrule
$n \le 15$         & enumerate every shape ($C_{15} = 9{,}694{,}845$ still fits) & \secref{c01-counting} \\
$n \le 10^{3}$     & $O(n^2)$ is fine -- the naive height-at-each-node passes  & \secref{c01-kinds} \\
$n \le 10^{6}$     & one visit per node, $O(1)$ combine -- single-pass recursion & \secref{c01-anatomy} \\
count ``modulo $10^9{+}7$'' & the answer is a huge formula -- Catalan/Cayley DP & \secref{c01-counting} \\
\bottomrule
\end{tabular}
\end{center}

\begin{notebox}[Not this chapter?]
If the tree must \emph{stay ordered} for fast lookup, you need BSTs (\chref{c05});
if it must \emph{support updates} while staying shallow, balanced trees (\chref{c06});
if queries ask about \emph{paths between two fixed nodes}, that is LCA (\chref{c11}).
This chapter supplies the vocabulary all of them speak.
\end{notebox}

% ---------------- Notation ----------------
\subsection*{Notation Used in This Chapter}

\begin{center}
\renewcommand{\arraystretch}{1.2}
\begin{tabular}{@{}ll@{\qquad}ll@{}}
\toprule
\textbf{Symbol} & \textbf{Meaning} & \textbf{Symbol} & \textbf{Meaning} \\
\midrule
$n$          & number of nodes               & $\depth(v)$   & edges from root down to $v$ \\
$T$          & a tree                        & $\height(v)$  & edges from $v$ to deepest leaf below \\
$\parent(v)$ & the node above $v$            & $\size(v)$    & nodes in the subtree rooted at $v$ \\
$\children(v)$ & nodes directly below $v$    & $\level(v)$   & same as $\depth(v)$ (root $=$ level $0$) \\
$\degop(v)$  & number of neighbours of $v$   & $\leaves(T)$  & the set of leaves of $T$ \\
$\subtree(v)$ & the tree under $v$           & $\nil$        & the empty / \code{null} child \\
$h$          & height of the whole tree      & $C_n$         & $n$-th Catalan number \\
\bottomrule
\end{tabular}
\end{center}

Throughout the book $\height$ and $\depth$ are counted in \textbf{edges}, so a single
lone node has height $0$ and $\height(\nil)=-1$. The convention note in Appendix~A
explains why this matters before you submit.
